\textbf{Phần 1. Trắc nghiệm lựa chọn}
\setcounterpageref{ex}{0}

\begin{ex} %[4D1V03]
Cho biểu thức $A = 1 \times 2 \times 3 \times \dots \times 9 + 2025$. Giá trị của biểu thức $A$ là
\choice
{Một số chẵn và chia hết cho 2}
{Một số có chữ số tận cùng bằng 0}
{\True Một số lẻ và không chia hết cho 2}
{Một số chẵn và chia hết cho 5}
\loigiai{
Tích $1 \times 2 \times 3 \times \dots \times 9$ có chứa các thừa số chẵn ($2, 4, 6, 8$) nên tích này là một số chẵn.\\
Số $2025$ có chữ số tận cùng là 5 nên là một số lẻ.\\
Ta có: $\text{Chẵn} + \text{Lẻ} = \text{Lẻ}$.\\
Vậy giá trị của biểu thức $A$ là một số lẻ và không chia hết cho 2.\\
\textbf{Đáp án của câu hỏi là C}
}
\end{ex}

\begin{ex} %[4D1V03]
Có bao nhiêu số chẵn có 4 chữ số khác nhau được lập từ các chữ số $0; 1; 2; 3; 4$?
\choice
{\True $60$ số}
{$48$ số}
{$72$ số}
{$64$ số}
\loigiai{
Số chẵn có 4 chữ số có dạng $\overline{abcd}$ với $d \in \{0; 2; 4\}$.\\
\textbf{Trường hợp 1:} $d = 0$.\\
- Chữ số $a$ có 4 cách chọn (từ $1; 2; 3; 4$).\\
- Chữ số $b$ có 3 cách chọn.\\
- Chữ số $c$ có 2 cách chọn.\\
Số các số lập được là: $4 \times 3 \times 2 = 24$ (số).\\
\textbf{Trường hợp 2:} $d \in \{2; 4\}$ (có 2 cách chọn $d$).\\
- Chữ số $a$ có 3 cách chọn (khác 0 và khác $d$).\\
- Chữ số $b$ có 3 cách chọn (khác $a$ và $d$, kể cả 0).\\
- Chữ số $c$ có 2 cách chọn.\\
Số các số lập được là: $2 \times 3 \times 3 \times 2 = 36$ (số).\\
Tổng số các số chẵn lập được là: $24 + 36 = 60$ (số).\\
\textbf{Đáp án của câu hỏi là A}
}
\end{ex}

\begin{ex} %[4D1V03]
Tổng của ba số tự nhiên liên tiếp bằng $84$. Tích của ba số tự nhiên đó là
\choice
{$21\,900$}
{$22\,450$}
{$21\,840$}
{\True $21\,924$}
\loigiai{
Vì tổng của ba số tự nhiên liên tiếp bằng $84$ nên số ở giữa (số thứ hai) là:
$$84 : 3 = 28$$
Ba số tự nhiên liên tiếp đó lần lượt là: $27; 28; 29$.\\
Tích của ba số đó là:
$$27 \times 28 \times 29 = 756 \times 29 = 21\,924$$
\textbf{Đáp án của câu hỏi là D}
}
\end{ex}

\begin{ex} %[4D1V03]
Giữa hai số chẵn có đúng $18$ số lẻ liên tiếp. Hiệu của hai số chẵn đó bằng
\choice
{$18$}
{\True $36$}
{$37$}
{$35$}
\loigiai{
Mỗi số lẻ liên tiếp cách nhau 2 đơn vị.\\
Giữa hai số chẵn có 18 số lẻ liên tiếp nên khoảng cách giữa hai số chẵn đó gồm $18$ khoảng cách 2 đơn vị.\\
Hiệu của hai số chẵn đó là:
$$18 \times 2 = 36$$
\textbf{Đáp án của câu hỏi là B}
}
\end{ex}

\begin{ex} %[4D1V03]
Tìm số lẻ bé nhất có 5 chữ số khác nhau mà tổng các chữ số của số đó bằng $20$.
\choice
{$10\,379$}
{$10\,469$}
{\True $10\,289$}
{$12\,089$}
\loigiai{
Để số là bé nhất có 5 chữ số khác nhau:\\
- Chữ số hàng chục nghìn bé nhất khác 0 là 1.\\
- Chữ số hàng nghìn bé nhất khác 1 là 0.\\
Tổng 3 chữ số còn lại là: $20 - (1 + 0) = 19$.\\
Vì số cần tìm là số lẻ nên chữ số hàng đơn vị phải là số lẻ. Để số bé nhất, chữ số hàng đơn vị nên lớn nhất có thể để các hàng trăm và hàng chục nhỏ lại:\\
- Chọn chữ số hàng đơn vị là 9.\\
Khi đó tổng chữ số hàng trăm và hàng chục là: $19 - 9 = 10$.\\
Để số bé nhất, chọn chữ số hàng trăm bé nhất khác 1, 0, 9 là 2; khi đó chữ số hàng chục là $10 - 2 = 8$.\\
Vậy số cần tìm là $10\,289$.\\
\textbf{Đáp án của câu hỏi là C}
}
\end{ex}

\begin{ex} %[4D1V03]
Tìm số lớn trong hai số lẻ có tổng bằng $2024$, biết giữa chúng có đúng $15$ số chẵn liên tiếp.
\choice
{\True $1027$}
{$1025$}
{$997$}
{$1032$}
\loigiai{
Giữa hai số lẻ có đúng 15 số chẵn liên tiếp thì hiệu giữa hai số lẻ đó là:
$$15 \times 2 = 30$$
Số lớn là:
$$(2024 + 30) : 2 = 1027$$
\textbf{Đáp án của câu hỏi là A}
}
\end{ex}

\begin{ex} %[4D1V03]
Cho dãy số gồm các số chẵn liên tiếp: $2; 4; 6; 8; \dots; 200$. Chữ số thứ $100$ của dãy số tính từ trái sang phải là chữ số nào?
\choice
{$1$}
{$0$}
{$4$}
{\True $2$}
\loigiai{
- Các số chẵn có 1 chữ số là: $2; 4; 6; 8$ (có 4 số, tương ứng có 4 chữ số).\\
- Số chữ số còn lại đến chữ số thứ 100 là: $100 - 4 = 96$ (chữ số).\\
- Các số chẵn có 2 chữ số từ 10 đến 98 có số lượng là:
$$(98 - 10) : 2 + 1 = 45\text{ (số)}$$
45 số này gồm: $45 \times 2 = 90$ (chữ số).\\
- Số chữ số còn lại là: $96 - 90 = 6$ (chữ số).\\
6 chữ số này thuộc các số chẵn có 3 chữ số bắt đầu từ 100:\\
Số chẵn 100 chiếm 3 chữ số (chữ số thứ 95, 96, 97).\\
Số chẵn 102 chiếm 3 chữ số (chữ số thứ 98, 99, 100).\\
Vậy chữ số thứ 100 của dãy số là chữ số $2$ (chữ số tận cùng của số 102).\\
\textbf{Đáp án của câu hỏi là D}
}
\end{ex}

\begin{ex} %[4D1V03]
Một giá sách có hai ngăn chứa tất cả $180$ quyển sách. Nếu chuyển $15$ quyển từ ngăn trên xuống ngăn dưới thì số sách ở ngăn trên vẫn nhiều hơn ngăn dưới $10$ quyển. Ban đầu ngăn trên có số quyển sách là
\choice
{$105$ quyển}
{\True $110$ quyển}
{$115$ quyển}
{$120$ quyển}
\loigiai{
Sau khi chuyển 15 quyển sách, ngăn trên nhiều hơn ngăn dưới 10 quyển.\\
Vậy ban đầu ngăn trên nhiều hơn ngăn dưới số quyển sách là:
$$10 + 15 \times 2 = 40\text{ (quyển)}$$
Ban đầu ngăn trên có số quyển sách là:
$$(180 + 40) : 2 = 110\text{ (quyển)}$$
\textbf{Đáp án của câu hỏi là B}
}
\end{ex}

\begin{ex} %[4D1V03]
Tổng của 20 số chẵn liên tiếp đầu tiên lớn hơn 0 là
\choice
{$400$}
{$410$}
{\True $420$}
{$440$}
\loigiai{
Dãy 20 số chẵn liên tiếp đầu tiên lớn hơn 0 là: $2; 4; 6; \dots; 40$.\\
Số số hạng là 20.\\
Tổng của 20 số chẵn đó là:
$$(2 + 40) \times 20 : 2 = 42 \times 10 = 420$$
\textbf{Đáp án của câu hỏi là C}
}
\end{ex}

\begin{ex} %[4D1V03]
Cho 5 chữ số $0; 2; 3; 5; 7$. Tổng của số lớn nhất có 5 chữ số khác nhau và số lẻ bé nhất có 5 chữ số khác nhau lập từ 5 chữ số trên là
\choice
{\True $95\,677$}
{$95\,657$}
{$96\,577$}
{$94\,877$}
\loigiai{
- Số lớn nhất có 5 chữ số khác nhau lập từ 5 chữ số trên là $75\,320$.\\
- Số lẻ bé nhất có 5 chữ số khác nhau lập từ 5 chữ số trên là $20\,357$ (chữ số hàng đơn vị lẻ là 7 để các hàng cao chọn bé nhất: $2; 0; 3; 5$).\\
Tổng của hai số đó là:
$$75\,320 + 20\,357 = 95\,677$$
\textbf{Đáp án của câu hỏi là A}
}
\end{ex}

\begin{ex} %[4D1V03]
Người ta viết các số chẵn liên tiếp thành dãy: $2; 4; 6; 8; 10; \dots$ Số chẵn thứ $25$ trong dãy số này là
\choice
{$48$}
{$52$}
{$46$}
{\True $50$}
\loigiai{
Công thức tìm số hạng thứ $n$ của dãy số chẵn bắt đầu từ 2 là:
$$\text{Số thứ } n = 2 + (n - 1) \times 2 = 2 \times n$$
Số chẵn thứ 25 trong dãy là:
$$2 \times 25 = 50$$
\textbf{Đáp án của câu hỏi là D}
}
\end{ex}

\begin{ex} %[4D1V03]
Có bao nhiêu số tự nhiên có 3 chữ số mà trong mỗi số đó tất cả các chữ số đều là số chẵn?
\choice
{$80$ số}
{\True $100$ số}
{$125$ số}
{$64$ số}
\loigiai{
Các chữ số chẵn gồm: $\{0; 2; 4; 6; 8\}$ (có 5 chữ số).\\
Số có 3 chữ số gồm toàn chữ số chẵn có dạng $\overline{abc}$:\\
- Chữ số $a$ có 4 cách chọn (chọn trong $\{2; 4; 6; 8\}$, vì $a \neq 0$).\\
- Chữ số $b$ có 5 cách chọn (chọn trong $\{0; 2; 4; 6; 8\}$).\\
- Chữ số $c$ có 5 cách chọn (chọn trong $\{0; 2; 4; 6; 8\}$).\\
Số lượng các số thỏa mãn là:
$$4 \times 5 \times 5 = 100\text{ (số)}$$
\textbf{Đáp án của câu hỏi là B}
}
\end{ex}

\textbf{Phần 2. Trắc nghiệm đúng sai}
\setcounterpageref{ex}{0}

\begin{ex} %[4D1V03]
Xét tính đúng sai của các khẳng định sau về số chẵn và số lẻ:
\choiceTF
{\True Tích của hai số chẵn liên tiếp luôn là một số chia hết cho 8.}
{Tổng của 3 số lẻ liên tiếp bất kì luôn là một số chẵn.}
{\True Giữa hai số lẻ liên tiếp không có số lẻ nào và có đúng 1 số chẵn.}
{\True Có tất cả 450 số chẵn có 3 chữ số.}
\loigiai{
\begin{itemize}
    \item Trong hai số chẵn liên tiếp, luôn có một số chia hết cho 2 và một số chia hết cho 4. Do đó tích của chúng luôn chia hết cho $2 \times 4 = 8$. Khẳng định này Đúng.
    \item Ta có $\text{Lẻ} + \text{Lẻ} + \text{Lẻ} = \text{Lẻ}$. Do đó tổng của 3 số lẻ liên tiếp luôn là một số lẻ (không phải số chẵn). Khẳng định này Sai.
    \item Hai số lẻ liên tiếp cách nhau 2 đơn vị, ở giữa chúng là số tự nhiên duy nhất và là số chẵn. Khẳng định này Đúng.
    \item Các số chẵn có 3 chữ số từ 100 đến 998 có: $(998 - 100) : 2 + 1 = 450$ (số). Khẳng định này Đúng.
\end{itemize}
\begin{itemize}
    \item a) \textbf{Đúng}.
    \item b) \textbf{Sai}.
    \item c) \textbf{Đúng}.
    \item d) \textbf{Đúng}.
\end{itemize}
}
\end{ex}

\begin{ex} %[4D1V03]
Cho một dãy gồm 5 số lẻ liên tiếp có tổng bằng $105$. Xét tính đúng sai của các khẳng định sau:
\choiceTF
{\True Số bé nhất trong 5 số đó là 17.}
{Số lớn nhất trong 5 số đó là 27.}
{\True Tích của số bé nhất và số lớn nhất trong dãy bằng 425.}
{\True Hiệu giữa số lớn nhất và số bé nhất trong dãy bằng 8.}
\loigiai{
Vì có 5 số lẻ liên tiếp nên số ở chính giữa (số thứ ba) là: $105 : 5 = 21$.\\
Năm số lẻ liên tiếp đó là: $17; 19; 21; 23; 25$.
\begin{itemize}
    \item Số bé nhất trong dãy là 17. Khẳng định a Đúng.
    \item Số lớn nhất trong dãy là 25 (không phải 27). Khẳng định b Sai.
    \item Tích số bé nhất và số lớn nhất: $17 \times 25 = 425$. Khẳng định c Đúng.
    \item Hiệu giữa số lớn nhất và bé nhất: $25 - 17 = 8$. Khẳng định d Đúng.
\end{itemize}
\begin{itemize}
    \item a) \textbf{Đúng}.
    \item b) \textbf{Sai}.
    \item c) \textbf{Đúng}.
    \item d) \textbf{Đúng}.
\end{itemize}
}
\end{ex}

\textbf{Phần 3. Trả lời ngắn}
\setcounterpageref{ex}{0}

\begin{ex} %[4D1V03]
Tìm số chẵn lớn nhất có 5 chữ số khác nhau mà chữ số hàng chục nghìn bằng $7$.
\shortans{79864}
\loigiai{
\textbf{Đáp số:} $79864$\\
Để số chẵn có 5 chữ số khác nhau lớn nhất với chữ số hàng chục nghìn là 7:\\
- Hàng nghìn chọn chữ số lớn nhất chưa dùng là 9.\\
- Hàng trăm chọn chữ số lớn nhất chưa dùng là 8.\\
- Hàng chục chọn chữ số lớn nhất chưa dùng là 6.\\
- Hàng đơn vị phải là chữ số chẵn lớn nhất chưa dùng khác 8, 6 nên chọn 4.\\
Số cần tìm là $79\,864$.
}
\end{ex}

\begin{ex} %[4D1V03]
Tìm số lớn trong hai số lẻ có tổng bằng $156$, biết giữa chúng có đúng $8$ số lẻ khác.
\shortans{87}
\loigiai{
\textbf{Đáp số:} $87$\\
Giữa hai số lẻ có 8 số lẻ khác thì hai số lẻ đó cách nhau:
$$(8 + 1) \times 2 = 18\text{ (đơn vị)}$$
Số lớn là:
$$(156 + 18) : 2 = 87$$
}
\end{ex}

\begin{ex} %[4D1V03]
Có bao nhiêu số tự nhiên có 3 chữ số chia hết cho $2$ nhưng không chia hết cho $5$?
\shortans{360}
\loigiai{
\textbf{Đáp số:} $360$\\
- Các số có 3 chữ số chia hết cho 2 là các số từ 100 đến 998:\\
$$(998 - 100) : 2 + 1 = 450\text{ (số)}$$
- Các số có 3 chữ số chia hết cho cả 2 và 5 là các số có tận cùng là 0 (từ 100 đến 990):\\
$$(990 - 100) : 10 + 1 = 90\text{ (số)}$$
Số các số có 3 chữ số chia hết cho 2 nhưng không chia hết cho 5 là:
$$450 - 90 = 360\text{ (số)}$$
}
\end{ex}

\begin{ex} %[4D1V03]
Tính tổng của tất cả các số lẻ có hai chữ số lớn hơn $30$.
\shortans{2275}
\loigiai{
\textbf{Đáp số:} $2275$\\
Dãy các số lẻ có hai chữ số lớn hơn 30 là: $31; 33; 35; \dots; 99$.\\
Số lượng các số hạng là:
$$(99 - 31) : 2 + 1 = 35\text{ (số)}$$
Tổng của dãy số là:
$$(31 + 99) \times 35 : 2 = 130 \times 35 : 2 = 2275$$
}
\end{ex}

\begin{ex} %[4D1V03]
Từ $1$ đến $100$, có bao nhiêu chữ số chẵn được dùng để viết các số trong dãy này?
\shortans{90}
\loigiai{
\textbf{Đáp số:} $90$\\
Các chữ số chẵn gồm: $0; 2; 4; 6; 8$.\\
- Ở hàng đơn vị: Cứ 10 số liên tiếp có đúng 5 chữ số chẵn ở hàng đơn vị. Từ 1 đến 100 có 10 nhóm 10 số nên số chữ số chẵn ở hàng đơn vị là:
$$10 \times 5 = 50\text{ (chữ số)}$$
- Ở hàng chục: Các số có chữ số hàng chục là số chẵn gồm các số trong các khoảng: từ 20 đến 29 (10 số), từ 40 đến 49 (10 số), từ 60 đến 69 (10 số), từ 80 đến 89 (10 số). Số chữ số chẵn ở hàng chục là:
$$4 \times 10 = 40\text{ (chữ số)}$$
- Số 100 có chữ số hàng trăm là 1 (lẻ), hai chữ số 0 đã được tính ở hàng chục và đơn vị.\\
Tổng số chữ số chẵn được dùng là:
$$50 + 40 = 90\text{ (chữ số)}$$
}
\end{ex}

\begin{ex} %[4D1V03]
Một cuốn sách được đánh số trang từ $1$ đến $120$. Hỏi có bao nhiêu trang sách được đánh số bằng một số lẻ?
\shortans{60}
\loigiai{
\textbf{Đáp số:} $60$\\
Các trang sách được đánh số bằng số lẻ tạo thành dãy số: $1; 3; 5; \dots; 119$.\\
Số trang sách được đánh số bằng số lẻ là:
$$(119 - 1) : 2 + 1 = 60\text{ (trang)}$$
}
\end{ex}
